% 2.10 — brief en documento/investigacion/2-10-inteligencia-de-negocios.md
\subsection{INTELIGENCIA DE NEGOCIOS}
\subsubsection{DEFINICIÓN Y REPORTES DESCRIPTIVOS}

La inteligencia de negocios (\textit{business intelligence}, BI) designa el conjunto de tecnologías de apoyo a la decisión que permiten a ejecutivos, gerentes y analistas tomar decisiones mejores y más rápidas \parencite{Chaudhuri2011}. \textcite{Negash2004} la formula desde los sistemas de información: los sistemas de BI combinan datos operativos con herramientas analíticas para presentar información compleja a quienes planifican y deciden, con el fin de mejorar la oportunidad y la calidad de los insumos de la decisión. \textcite{Watson2007} la describen como un proceso con dos actividades principales: introducir los datos (\textit{getting data in}), que corresponde al almacenamiento de datos tradicional y consiste en mover datos desde los sistemas fuente hacia un almacén integrado, y extraerlos (\textit{getting data out}), cuando usuarios y aplicaciones acceden a ese almacén para reportes empresariales, consultas, procesamiento analítico y analítica predictiva. Los mismos autores advierten que la primera actividad aporta por sí sola un valor limitado: la organización solo obtiene el valor completo cuando usa los datos para decidir.

Esa arquitectura tiene componentes reconocibles. \textcite{Chaudhuri1997} presentan el almacén de datos y el procesamiento analítico en línea (\textit{online analytical processing}, OLAP) como elementos esenciales del apoyo a la decisión y recorren sus piezas: herramientas de extracción, limpieza y carga (los procesos ETL de la sección 2.3), modelos de datos multidimensionales organizados en medidas y dimensiones, herramientas cliente de consulta y análisis, y gestión de metadatos. \textcite{Kimball2013} sistematizan el diseño de esa capa mediante el modelado dimensional: tablas de hechos que registran las medidas del negocio y tablas de dimensión que las describen, dispuestas en esquema en estrella. La aplicación de presentación con la que interactúa el usuario es el último eslabón de una cadena cuya calidad depende de las capas anteriores \parencite{Chaudhuri2011}.

El producto más visible de la BI es el reporte descriptivo. \textcite{Lepenioti2020} ordenan la analítica de negocios en tres niveles: la analítica descriptiva, que muestra lo que ocurrió; la predictiva, que estima lo que ocurrirá; y la prescriptiva, que busca el mejor curso de acción futuro y constituye, según los autores, el siguiente paso en la madurez analítica, mientras que el foco académico e industrial sigue puesto en las dos primeras. Los reportes descriptivos suelen entregarse como tableros (\textit{dashboards}): \textcite{Eckerson2010} los caracteriza como una capa de presentación sobre la infraestructura de BI que integra el monitoreo, el análisis y la gestión de indicadores de desempeño, y \textcite{Few2013} aporta los criterios de diseño visual para que la información más importante pueda monitorearse de un vistazo en una sola pantalla. La BI de autoservicio, por último, aspira a que usuarios no especialistas exploren los datos y deriven información accionable sin depender de los especialistas de BI \parencite{Alpar2016}.

Estas nociones fijan el lugar de la BI en el proyecto. El pipeline de datos de la sección 2.3 es la actividad de introducir los datos; los reportes de Power BI y el tablero del ERP en Angular son dos canales para extraerlos. El marco evolutivo de \textcite{Chen2012}, que distingue etapas de la BI y la analítica según los datos que procesan, ubica a este trabajo en su punto de partida: datos estructurados de un sistema transaccional. La separación entre canales es una decisión de arquitectura y no una definición, puesto que \textcite{Watson2007} incluyen la analítica predictiva dentro de la BI\@. Con la tríada de \textcite{Lepenioti2020} como criterio, se prevé que Power BI cubra la analítica descriptiva del ERP (horas, compras, presupuesto ejecutado y certificaciones por centro de costo), que el módulo de aprendizaje automático cubra la predictiva dentro de la interfaz transaccional, y que la prescriptiva quede como trabajo futuro. Las razones técnicas de esa separación se desarrollan en la subsección siguiente.

\subsubsection{POWER BI}

Power BI es la plataforma de analítica de negocios de Microsoft, un conjunto integrado de herramientas para conectar, visualizar y compartir datos, con dos componentes principales: Power BI Desktop, aplicación gratuita para el modelado y la creación de reportes, y el servicio Power BI, orientado a compartir mediante espacios de trabajo y refresco programado \parencite{MSPowerBIOverview, MSDAXOverview}. La preparación de los datos corre por cuenta de Power Query, un motor que realiza la extracción, transformación y carga (ETL) desde cientos de fuentes, registra cada transformación como un paso de consulta en el lenguaje M y la aplica sin modificar la fuente \parencite{MSPowerQueryWhatIs}. El modelado se apoya en DAX (\textit{Data Analysis Expressions}), lenguaje de fórmulas sobre tablas relacionadas de un modelo tabular con el que se definen medidas, columnas calculadas y funciones de inteligencia de tiempo \parencite{MSDAXOverview}. El modelo semántico (antes \textit{dataset}) es el contenedor sobre el que se construyen los reportes, y \textit{dashboard} designa allí una pantalla de mosaicos \parencite{MSPowerBIServiceConcepts}; en este documento se reserva ``tablero'' para el sentido de gestión de la subsección anterior.

El modo de acceso a la fuente condiciona el diseño. En el modo de importación los datos se cargan en la caché en memoria del modelo semántico y los visuales no reflejan cambios en la fuente hasta el siguiente refresco; en DirectQuery cada visual consulta la fuente en el momento. Microsoft recomienda importar por defecto, y en su documentación las fuentes habilitadas para DirectQuery son bases de datos, mientras que los sitios web figuran entre las fuentes de importación \parencite{MSPowerBIDirectQuery}. Importar implica refrescar: en capacidad compartida el servicio limita cada modelo a ocho refrescos programados por día \parencite{MSPowerBIRefresh}. Si el servicio no puede establecer una conexión de red directa con la fuente, el refresco requiere una puerta de enlace de datos local (\textit{on-premises data gateway}) instalada en la red de la empresa \parencite{MSPowerBIRefresh, MSPowerBIGateway}; como la fase 1 corre en un servidor propio, su necesidad depende de si la API del ERP es alcanzable desde Internet, condición que el Capítulo III debe resolver. Estas propiedades fundamentan la decisión anterior: una copia importada que se actualiza a lo sumo ocho veces al día sirve para reportes descriptivos, no para predicciones calculadas bajo demanda sobre proyectos activos y atadas a la sesión del usuario.

El segundo condicionante es la autenticación. Power BI llega a la API de NestJS por el conector Web, cuya función \texttt{Web.Contents} admite, entre otros, los tipos Anonymous, Web API (clave de API) y Organizational Account, permite definir encabezados HTTP y envuelve una respuesta JSON con \texttt{Json.Document}; las solicitudes POST solo pueden hacerse de forma anónima \parencite{MSPowerQueryWeb}. El flujo OAuth integrado de Organizational Account solo se inicia cuando el servicio responde 401 con un encabezado que remite a Microsoft Entra ID; más control sobre el flujo exige un conector personalizado \parencite{MSPowerQueryAuth}, que puede declarar el tipo OAuth e implementar \texttt{StartLogin}, \texttt{FinishLogin} y \texttt{Refresh} con PKCE, o el tipo Key, cuya clave el motor inserta en el encabezado \texttt{Authorization} \parencite{MSPowerQueryHandlingAuth}. Como el proveedor de identidad del proyecto es Keycloak y no Entra ID, quedan dos caminos abiertos: un conector personalizado con OAuth y PKCE contra Keycloak, o una credencial de servicio emitida por Keycloak cuyo token se pasa en el encabezado desde \texttt{Web.Contents}; la restricción sobre POST afecta a la segunda si el intercambio de credenciales se hiciera desde M\@. La elección se justifica en el Capítulo III y se construye en el incremento 5.

La elección frente a otras herramientas responde a un criterio técnico derivado del alcance: la BI debe leer el ERP por una API protegida, no por acceso directo a PostgreSQL\@. Apache Superset se conecta a bases de datos SQL \parencite{ApacheSuperset}; Metabase documenta controladores para bases de datos y no menciona API REST como fuente \parencite{MetabaseDatabases}; Tableau accede a fuentes HTTP sin conector nativo mediante un \textit{web data connector} en JavaScript, cuya versión 2.0 está obsoleta desde la versión 2023.1 \parencite{TableauWDC}. Power BI trae de fábrica el conector Web con JSON y encabezados, y su herramienta de autoría es gratuita. El cuadro~\ref{tab:bi-alternativas} resume la comparación; no se afirma una superioridad general, sino la adecuación a este caso. Con esa base, se prevé publicar el modelo semántico en el servicio Power BI y compartir los reportes con los jefes de proyecto y administración en la modalidad de autoservicio descrita antes.

\begin{table}[H]
  \centering
  \caption{Herramientas de BI consideradas y su forma de acceso a los datos del ERP}\label{tab:bi-alternativas}
  \begin{tabular}{@{}>{\RaggedRight\arraybackslash}p{0.16\textwidth}>{\RaggedRight\arraybackslash}p{0.36\textwidth}>{\RaggedRight\arraybackslash}p{0.42\textwidth}@{}}
    \toprule
    \textbf{Herramienta} & \textbf{Fuentes documentadas} & \textbf{Acceso a la API del ERP} \\
    \midrule
    Power BI & Más de 100 fuentes: bases de datos, servicios en la nube, archivos y fuentes web & Conector Web de fábrica: JSON, encabezados HTTP, clave de API\@. OAuth integrado solo con Entra ID; otro proveedor exige conector personalizado \\
    Apache Superset & Bases de datos SQL & Sin conector REST documentado; requeriría acceso directo a la base, excluido por el alcance \\
    Metabase & Controladores para bases de datos (PostgreSQL, MySQL, SQL Server, entre otras) & Sin conector REST documentado; misma restricción \\
    Tableau & Conectores nativos y \textit{web data connector} para fuentes HTTP & Requiere programar un conector en JavaScript (WDC 3.0; la 2.0 está obsoleta) \\
    \bottomrule
  \end{tabular}
  \fuente{Elaboración propia a partir de \textcite{MSPowerBIOverview}, \textcite{MSPowerQueryAuth}, \textcite{ApacheSuperset}, \textcite{MetabaseDatabases} y \textcite{TableauWDC}.}
\end{table}
